\SourcePage{590}{593}
\section{Nonspherical nuclei}

A system of particles moving in a spherically symmetric field cannot possess a
rotational energy spectrum; indeed, in quantum mechanics the very notion of
rotation has no meaning for such a system. This applies also to the shell model
of the nucleus considered in the preceding section, with its spherically
symmetric self-consistent field.

In quantum mechanics, a rigorous division of a system's energy into internal
and rotational parts has no mean\SourcePageJoin{591}{594}ing in general. Such a division can only be
approximate, and is possible when physical circumstances make it a good
approximation to regard the system as a collection of particles moving in a
prescribed field that lacks spherical symmetry. The rotational structure of
the levels then arises when allowance is made for this field to rotate relative
to a fixed coordinate system. We encountered such a case, for example, in
molecules, whose electronic terms may be defined as the energy levels of the
electron system moving in the prescribed field of fixed nuclei.

Experiment shows that most nuclei do not, in fact, have a rotational
structure. For them, therefore, a spherically symmetric self-consistent field
is a good approximation: apart from quantum fluctuations, these nuclei are
spherical.

There is, however, a class of nuclei whose energy spectrum is of rotational
type (this class includes nuclei in the approximate atomic-weight ranges
$150<A<190$ and $A>220$). This property means that the approximation of a
spherically symmetric self-consistent field is wholly unsuitable for them. In
principle, their self-consistent field must be sought without any prior
assumption about its symmetries, so that the shape of the nucleus is itself
determined ``self-consistently.'' Experiment indicates that the appropriate
model for nuclei in this class is a self-consistent field having an axis of
symmetry and a symmetry plane perpendicular to it (that is, the symmetry of an
ellipsoid of revolution). The conception of nonspherical nuclei was developed
most fully in the work of \emph{A.~Bohr and B.~R.~Mottelson} (1952--1953).

We emphasize that these are two qualitatively different classes of nuclei.
One manifestation of this distinction is that a nucleus is either spherical
or nonspherical with a degree of nonsphericity that is by no means small.

Unfilled shells in a nucleus promote the appearance of nonsphericity; nucleon
pairing also appears to play an important part in this phenomenon. Closed
shells, by contrast, favor a spherical nucleus. A characteristic example is
the doubly magic nucleus ${}^{208}_{82}\mathrm{Pb}$. Because its nucleon
configuration is very strongly closed, this nucleus (and nuclei near it) is
spherical, thereby producing the interruption in the sequence of nonspherical
heavy nuclei.

\SourcePage{592}{595}
The energy levels of a nonspherical nucleus are represented as the sum of two
parts: the levels of the ``stationary'' nucleus and the energy of its rotation
as a whole. In even--even nuclei, the intervals within this rotational
structure are small compared with the separations between levels of the
``stationary'' nucleus.

The classification of the levels of a nonspherical nucleus is in many respects
analogous to the classification of the levels of a diatomic molecule composed
of identical atoms, since in the two cases the fields in which the particles
(nucleons or electrons) move have the same symmetry. We may therefore use
directly a number of the results obtained in Chapter~XI\footnote{We stress that
the analogy concerns the classification of the levels of a diatomic molecule,
not of a symmetric top. For particles moving in an axially symmetric field,
rotation about the field axis has no meaning, just as rotation about any axis
has no meaning for a system in a centrally symmetric field.}.

Let us first consider the classification of states of the ``stationary
nucleus.'' In an axially symmetric field, only the projection of the angular
momentum on the symmetry axis is conserved. Each nuclear state is consequently
characterized first of all by the magnitude $\Omega$ of the projection of its
total angular momentum\footnote{By definition, $\Omega\geqslant0$ (as with
the positive quantum number $\Lambda$ in diatomic molecules). Recall that
negative values of $\Omega$ for diatomic molecules could occur only because
$\Omega$ was defined as the sum $\Lambda+\Sigma$, while $\Sigma$ could be
positive or negative according to the relative directions of orbital angular
momentum and spin.}; this magnitude may be integral or half-integral. According
to the behavior of the wave function when the signs of the coordinates of all
nucleons relative to the nuclear center are reversed, the levels are divided
into even ($g$) and odd ($u$) levels.

For $\Omega=0$ a further distinction is made between positive and negative
states, according to the behavior of the wave function under reflection in a
plane through the nuclear axis (see \S~78).

The ground states of nonspherical even--even nuclei are $0_g$ states (the
numeral gives the value of $\Omega$), corresponding to zero angular momentum
and to the highest symmetry of the wave function. This is a consequence of the
pairwise pairing of all the neutrons and all the protons. If, however, a
nucleus contains an odd number of protons or neutrons, one may consider the
state of the ``odd'' nucleon in the self-consistent field of the even--even
nuclear ``core.''

The value of $\Omega$ is then determined by the projection of this nucleon's
angular momentum. Similarly, in an odd--odd nucleus $\Omega$ is obtained from
the angular-momentum projections of the odd neutron and odd
\SourcePageJoin{593}{596}proton ($\Omega=|\omega_p\pm\omega_n|$).

At the same time, it must be stressed that definite values cannot be assigned
to the projections of the nucleon's orbital angular momentum and spin.

Although the spin--orbit coupling of a nucleon is small compared with the
energy of its interaction with the self-consistent field of the core, it is
not small compared with the spacings between neighboring energy levels of the
nucleon in that field. Yet precisely this latter condition would be required
for perturbation theory to apply and permit the nucleon's orbital angular
momentum and spin to be treated separately to a good approximation\footnote{In
spherical nuclei it was nevertheless possible to determine the value of $l$
by using conservation of parity together with conservation of angular
momentum.}.

We now turn to the rotational structure of the levels of a nonspherical
nucleus. Its intervals are small compared with the spin--orbit interaction of
the nucleons in the nucleus; this situation corresponds to case $a$ in the
theory of diatomic molecules (see \S~83).

The total angular momentum $\mathbf J$ of the rotating nucleus is, of course,
conserved. For a specified $\Omega$, its magnitude $J$ takes the sequence of
values beginning with $\Omega$:
\begin{equation}
 J=\Omega,\quad \Omega+1,\quad \Omega+2,\ldots
 \tag{119.1}
\end{equation}
(see (83.2)). There is an additional restriction on the possible values of
$J$ for nuclei with $\Omega=0$: in the states $0_g^+$ and $0_u^+$, $J$ takes
only even values, whereas in $0_g^-$ and $0_u^-$ it takes odd values (see
\S~86). In particular, for the rotational levels of the ground term of an
even--even nucleus ($0_g^+$), $J$ takes the values $0,2,4,\ldots$

The rotational energy of the nucleus is
\begin{equation}
 E_{\mathrm{rot}}=\frac{\hbar^2}{2I}J(J+1),
 \tag{119.2}
\end{equation}
where $I$ is the moment of inertia of the nucleus about an axis perpendicular
to its symmetry axis. This formula corresponds to the analogous expression in
the theory of diatomic molecules (the term depending

\SourcePage{594}{597}
on $J$ in (83.6)). The lowest level corresponds to the least possible value of
$J$, namely $J=\Omega$.

By (119.2), the rotational level structure obeys definite interval rules that,
for a specified $\Omega$, are independent of the other characteristics of the
level. Thus the components of the rotational structure of the ground term of
an even--even nucleus (with $J=2,4,6,8,\ldots$) lie above the lowest level
($J=0$) by intervals whose ratios are $1:3$, $3:7:12\ldots$

Formula (119.2), however, is insufficient for states with $\Omega=1/2$, which
may occur in nuclei containing an odd number of nucleons.

In this case the interaction of the odd nucleon with the centrifugal field of
the rotating nucleus makes a contribution to the energy comparable with
(119.2). Its dependence on $J$ can be found as follows.

As is known from mechanics (see Volume~I, \S~39), the energy of a particle in
a rotating coordinate system contains an additional term equal to the product
of the angular velocity and the particle's angular momentum. The corresponding
term in the nuclear Hamiltonian may be written as
$2b\hat{\mathbf K}\hat{\boldsymbol\sigma}$, where $b$ is a constant,
$\hat{\mathbf K}$ is the rotational angular momentum of the nuclear core (the
nucleus without the last nucleon), and $\hat{\boldsymbol\sigma}$ is the
angular momentum of the nucleon. The latter is to be understood here in a
purely formal sense (in reality, an angular-momentum vector of a nucleon in an
axial nuclear field does not exist): it is an operator analogous to the
spin-$1/2$ operator, producing transitions between states whose
angular-momentum projections are $\pm1/2$, in accordance with
$\Omega=1/2$\footnote{The special feature of the case $\Omega=1/2$ is
precisely the existence of matrix elements of the energy perturbation for
transitions between states that differ only in the sign of the
angular-momentum projection and therefore have the same energy. An energy
shift consequently appears already in first-order perturbation theory. This
phenomenon is analogous to the $\Lambda$-doubling of the levels of a diatomic
molecule with $\Omega=1/2$ (see \S~88).}. Since
$\mathbf K=\mathbf J-\boldsymbol\sigma$, the eigenvalues of this operator are
\[
 2b\mathbf K\boldsymbol\sigma
 =b\left[J(J+1)-K(K+1)-\frac34\right].
\]
Adding, for convenience, the $J$-independent constant $b/2$, we find this
quantity to be $\pm b(J+1/2)$ when $J=K\pm1/2$.

This expression can be written $(-1)^{J-1/2}b(J+1/2)$ on noting that the
angular momentum $K$ of the core, which is an even--even nucleus, is an even
integer. Thus,

\SourcePage{595}{598}
the final expression for the rotational energy of a nucleus with
$\Omega=1/2$ is
\begin{equation}
 E_{\mathrm{rot}}=\frac{\hbar^2}{2I}J(J+1)
 +(-1)^{J-1/2}b\left(J+\frac12\right)
 \tag{119.3}
\end{equation}
(A.~Bohr and B.~Mottelson, 1953). Notice that if $b$ is positive and
sufficiently large, the level with $J=3/2$ can lie below that with $J=1/2$.
The normal ordering of rotational levels, in which the lowest level has the
least possible $J$, may therefore be violated.

The moment of inertia of a nonspherical nucleus cannot be calculated as that
of a rigid body of prescribed shape. Such a calculation would be possible only
if the nucleons moving in the nuclear self-consistent field could be regarded
as having no direct interactions with one another. In reality, pairing reduces
the moment of inertia below the value for a rigid body.

The magnetic moment $\boldsymbol\mu$ of a nonspherical nucleus consists of the
magnetic moment of the ``stationary'' nucleus and a moment associated with the
rotation of the nucleus. After averaging over the motion of the nucleons in
the nucleus, the first is directed along the nuclear axis. Denoting its
magnitude by $\mu'$ and the unit vector along that axis by $\mathbf n$, we
write it as $\mu'\mathbf n$. The magnetic moment due to rotation, after the
same averaging, is directed along $\mathbf J-\Omega\mathbf n$, the total
mechanical angular momentum of the nucleus less the angular momentum of the
nucleons in the ``stationary nucleus''\footnote{This form can be used only
when $\Omega\ne1/2$ (see Problem~2).}.

Thus
\begin{equation}
 \boldsymbol\mu=\mu'\mathbf n+g_r(\mathbf J-\Omega\mathbf n).
 \tag{119.4}
\end{equation}
Here $g_r$ is the gyromagnetic factor for nuclear rotation. Since only the
protons contribute to the magnetic moment arising from rotation,
\begin{equation}
 g_r=\frac{I_p}{I_p+I_n},
 \tag{119.5}
\end{equation}
where $I_n$ and $I_p$ are the neutron and proton parts of the nuclear moment
of inertia (for a system consisting solely of protons one would simply have
$g_r=1$). In general, the ratio (119.5) does not equal the ratio $Z/A$ of the
number of protons to the total nuclear mass.

\SourcePage{596}{599}
After averaging over the nuclear rotation, the magnetic moment is directed
along the conserved vector $\mathbf J$:
\[
 \hat{\overline{\boldsymbol\mu}}
 =\frac\mu J\hat{\mathbf J}
 =(\mu'-\Omega g_r)\overline{\mathbf n}+g_r\hat{\mathbf J}.
\]
As usual, we multiply both sides of this equality by $\hat{\mathbf J}$ and
pass to the eigenvalues. In the nuclear ground state $\Omega=J$, and we obtain
\begin{equation}
 \mu=(\mu'+g_r)\frac{J}{J+1}.
 \tag{119.6}
\end{equation}

\ProblemHeading 1. Express the quadrupole moment $Q$ of a rotating nucleus in terms
of the quadrupole moment $Q_0$ relative to axes fixed in the nucleus (A.~Bohr,
1951).

\SolutionHeading The quadrupole-moment tensor operator of the rotating nucleus
is expressed in terms of $Q_0$ as
\[
 Q_{ik}=\frac32Q_0\left(n_in_k-\frac13\delta_{ik}\right);
\]
this is a symmetric tensor of zero trace, constructed from the components of
the unit vector $\mathbf n$ along the nuclear axis, with $Q_{zz}=Q_0$.
Averaging over the rotational state of the nucleus proceeds as in the solution
of the problem to \S~29, except that $n_iJ_i=\Omega$ rather than zero, and
gives an expression of the form (75.2), with
\[
 Q=Q_0\frac{3\Omega^2-J(J+1)}{(2J+3)(J+1)}.
\]
For the nuclear ground state with $\Omega=J$, this becomes
\[
 Q=Q_0\frac{(2J-1)J}{(2J+3)(J+1)}.
\]
As $J$ increases, the ratio $Q/Q_0$ tends to unity, though rather slowly.

\ProblemHeading 2. Determine the magnetic moment in the ground state of a nucleus
with $\Omega=1/2$.

\SolutionHeading In this case the magnetic-moment operator may be written,
using the operator $\hat{\boldsymbol\sigma}$ introduced in the text, as
\[
 \hat{\boldsymbol\mu}=2\mu'\hat{\boldsymbol\sigma}+g_r\hat{\mathbf K},
 \qquad \hat{\mathbf K}=\hat{\mathbf J}-\hat{\boldsymbol\sigma}.
\]
The remaining calculation is analogous to that carried out in the text. If
the nuclear ground level has $J=1/2$ (so that $K=J-1/2=0$), the result is
$\mu=\mu'$. If instead the ground state has $J=3/2$ (so that
$K=J+1/2=2$), then $\mu=\frac95g_r-\frac35\mu'$.

\ProblemHeading 3. Determine the energies of the first few levels in the rotational
structure of the ground state of an even--even nucleus having the symmetry of
a triaxial ellipsoid.

\SourcePage{597}{600}
\SolutionHeading The ground state of an even--even nucleus corresponds to the
most symmetric wave function of the ``stationary'' nucleus, namely a function
with the symmetry belonging to the representation $A_1$ of the group $D_2$.
For a specified $J$, there are therefore only $(J/2)+1$ distinct levels when
$J$ is even, or $(J-1)/2$ when $J$ is odd. For $J=2$ they are given by formula
(7) obtained in the problems to \S~103, and for $J=3$ by formula (8).
